\section{Background}
\label{sec:background}

\subsection{beni in brief}
\label{sec:bg-beni}

beni is a statically typed, strict, pure functional language in the family of Elm. Its compiler,
written in Zig, emits modern JavaScript as ES modules, and its primary platform is the browser. It
keeps Elm's guarantees---no run-time exceptions in ordinary code, exhaustive pattern matching,
managed effects---and departs from Elm in a few deliberate places. Those that matter here are the
following.

\paragraph{Syntax.} beni uses Unicode notation: \code{→} for arrows, \code{λ} for lambdas,
\code{▷} for the forward pipe, \code{…} for a spread, \code{⊤} for the unit type and value, and
\code{×} for tuple types. A function body is a block of bindings followed by its result. Lists are
written, built and matched with brackets and spreads: \code{[ x, …rest ]} matches a non-empty
list, \code{[ …xs, …ys ]} concatenates. Records are immutable and updated by \code{\{ r | f = v \}}.
A small \code{update} function in the style of a TodoMVC implementation looks like this:
\begin{Verbatim}
update : Msg, Model → Model
update msg model =
    case msg of
        Typed text →
            { model | draft = text }
        Add title →
            { model | todos = model.todos ++ [ { id = model.nextId, title = title, completed = False } ]
                    , nextId = model.nextId + 1 }
\end{Verbatim}

\paragraph{No automatic currying.} Function types are $n$-ary---\code{Msg, Model → Model} takes two
arguments---and every call is saturated. Partial application is written with a placeholder,
\code{f a \_}, which is sugar for the lambda \code{λp → f a p}; the pipe \code{x ▷ f a} is rewritten
to \code{f x a} before anything runs. There are no point-free composition operators. Every closure
in a beni program is therefore a lambda the programmer wrote (or a placeholder that desugars to
one), and the compiler can see what it captures.

\paragraph{Evaluation order.} beni is strict and evaluates left to right in source order: in an
application, the callee and then the arguments, left to right, and then the call; record literals
field by field in written order; a record update evaluates the record and then the updated fields;
block bindings in the order written; a self tail call evaluates all its new arguments before any
parameter is rebound. The order is part of the language definition and is observable today through
\code{Debug.log}; the compiler's release optimiser is constrained not to change it.

\paragraph{Inferred effects.} beni has no effect syntax. Its type checker infers, for every function
type, two bits---whether a call may perform an effect, and whether it may suspend the current
fiber---and solves them as a directed graph of inclusion edges (``this class is at least as
effectful as that one'') after each module is type-checked. The solved bits are published in the
module's interface and hashed with it. A function that takes a callback is summarised
\emph{polymorphically} in the callback's bits, so that \code{List.map} serves a pure and a
suspending callback with one definition. Concurrency is cooperative: fibers interleave only at
suspension points, and continuations are one-shot.

\paragraph{Static dispatch.} A type's methods are the public values of the module that declares
it; \code{x.m a} calls one, and a top-level annotation may carry a constraint such as
\code{where a.compare : a, a → Order}. The operators \code{==} and \code{<} call the receiver type's
\code{eq} and \code{compare}, derived structurally when the type declares none. Evidence for a
constraint is passed as a hidden argument.

\paragraph{The boundary.} Only privileged platform packages and the embedded core library may
declare a \code{foreign} value, which binds to an export of a sibling JavaScript file. Every
\code{foreign} already declares its effect rung (pure, impure or suspending), and the compiler
checks the sibling's export shape and arity.

\paragraph{The Elm Architecture.} A browser program is written in The Elm Architecture (TEA): a
\code{Model} type holding the application state, an \code{init} value, an
\code{update : Msg, Model → Model} function (or, with commands, one returning
\code{Model × Cmd Msg}) and a \code{view} from model to markup. A runtime owned by the platform
calls \code{update} for each message, replaces the current model with the result, and brings the
page up to date. \code{update} is where nearly all list edits in an application happen, and it is
where the runtime and the program both hold the model. Commands are thunks the runtime stores and
runs later; subscriptions are re-evaluated after each update.

beni has two browser back ends that matter here. The current one compiles each markup site into a
cloned template and patches dynamic holes with a reference check per hole, in the style of
SolidJS's compiler; it keeps each list it rendered so that it can skip the next walk when the list
is identical. A second, experimental back end---which we call the \emph{direct} platform---compiles
each message handler into a JavaScript function that runs the corresponding \code{update} arm and
writes the DOM directly, guided by a whole-program \emph{write-set analysis} that computes, for each
message, which paths of the model the arm may write. It keeps no rendered view: slots hold leaf
values, and per-row instance arrays hold the row items only to compare them by identity. That
design choice is what makes the hand-over protocol of \S\ref{sec:hard-tea} possible.

\subsection{How lists are represented today}
\label{sec:bg-lists}

\code{List} is beni's only sequence type. A list value is one of three JavaScript forms: a plain
array; a \emph{view}, an object \code{\{ b, o, length \}} giving an offset into a base array (used for
the \code{rest} of a pattern, in $O(1)$); or a 32-way radix trie with claimable head and tail
buffers. Every write is persistent: \code{List.set} copies a plain array of up to 256 elements and
path-copies a trie above that; \code{List.push} copies short arrays and claims the trie's tail
otherwise; prepending through \code{[ x, …xs ]} claims a free head slot. Readers outside the
library use three facts---\code{length}, \code{Array.isArray} and a \code{\$plain()} method that
materialises a plain array---so that the trie is invisible to them.

This design gives every program persistence at a cost that is modest per operation but paid
everywhere: the trie's code ships with every program that edits a list, every generic read path
can meet a non-plain form, and a list handed to JavaScript must first be flattened. A hand-written
JavaScript program would simply write \code{a[i] = v}.

\subsection{What in-place writes buy: a prototype study}
\label{sec:bg-measure}

Before designing a checker we measured what in-place writes would buy, in a prototype study that
transformed beni's own compiled output by hand, exactly as a compiler pass would, and ran six array
scenarios in Node~24 (five rounds; Chrome as a spot check). The baseline was a persistent
\emph{adaptive} array (plain copies up to 256 elements, a trie above). One of the variants, a purely
static one, wrote in place wherever last use and per-function interface summaries proved a value
unique; it is the run-time half of the checker this paper proposes. Its results relative to the
baseline:

\begin{itemize}
\item \emph{Building in a fold} (collecting by \code{push}, a histogram over 1\,000 or 100\,000
buckets, a coin-change table): 0.14--0.47 times the baseline's time, the largest gain in the study,
because every write is to an accumulator nobody else holds.
\item \emph{A grid game}, steady-state ticks of 1 or 100 cells on a board of one million cells:
0.44--0.48 times; the first tick, which must copy a board that the harness kept, was slower
(1.52 times), because a whole copy costs more than converting to a trie.
\item \emph{A TEA table} of 10\,000 rows: with the runtime keeping the rendered rows, as beni's
current runtime does, no counting or static scheme could write the rows in place---every write
found the rows shared. With a runtime that hands the model over, the writes landed in place, and
bought little that was visible (0.84--0.95 times on update and swap messages), because one in-place
element write saves about a microsecond next to a 40--60\,\textmu s render walk.
\item \emph{An undo history} that keeps every version gains nothing, correctly: nothing is unique.
\item \emph{Size}: the list runtime with the trie compressed to 1\,514 bytes (brotli); without it,
706 bytes.
\end{itemize}

The same study found that a static scheme that \emph{copies} where it cannot prove uniqueness is
harmless at a message boundary and catastrophic inside a loop: applied to a callback inside a fold,
it copied on every iteration and ran 8.2 times slower than the baseline. That observation is one
reason we propose an error rather than an optimisation: a silent copy makes the cost model depend
on facts the programmer cannot see, and in the worst case turns a linear loop quadratic.

\subsection{The question, stated precisely}
\label{sec:bg-question}

\emph{Editing} means a call of one of the core library's list-update primitives: \code{set},
\code{update}, \code{push}, \code{pop}, \code{swap}, \code{insertAt}, \code{removeAt}, and the two
the syntax lowers to, \code{cons} (\code{[ x, …xs ]}) and \code{append} (\code{[ …xs, …ys ]}, and
\code{++} on lists). \emph{Anything else} is any holder of the array, of any of the kinds listed in
\S\ref{sec:model-holders}, including the operand's own binding if it is read again after the edit.
\emph{Still holds} means the holder is live: some execution continuing from the edit reads through
it. \emph{The old version} is the very array, not an equal list. \emph{A compile error} means
neither a copy nor a warning. The one escape hatch is a new library function, \code{List.copy},
which makes intentional sharing explicit.

\paragraph{What ``observe'' means.} An in-place write is observable when some computation reads
the array's elements after the write through a reference it obtained before. Three kinds of reader
exist in beni: program code (\code{List.get}, a pattern, a fold, \code{==}, \code{Debug.toString}, a
markup hole that shows the list); the runtime (a renderer that keeps the last list it rendered); and
JavaScript (a sibling file that kept the array). Each of these is a \emph{content} holder. An
\emph{identity} holder keeps a reference only to ask \code{===} later. An in-place write does not
change identity, so an identity holder never sees a different object; what it may do is conclude
``unchanged'' from ``same object'', which is wrong after an in-place write. beni programs cannot
compare references, so this hazard belongs to the renderer, and we discharge it by an assumption on
the platforms (A4 in \S\ref{sec:soundness-assumptions}).

\paragraph{Two semantics.} The precise form of the guarantee is the one O'Connor, Linares
Ar\'evalo and Rizkallah give for uniqueness types: they work ``by statically ruling out every
program where the difference between the immutable and the mutable interpretations can be
observed'' \citep{oconnor2026separation}. Aspinall, Hofmann and Kone\v{c}n\'y prove such a theorem
for a first-order language \citep[Thm.~5.8]{aspinall2008usage}, and Linear Haskell proves one for
its mutable arrays \citep{bernardy2018linear}. We define the \emph{value semantics} $S_v$---every
update builds a new array, as beni does today---and the \emph{update semantics} $S_u$---every
accepted update writes its operand's array---and claim that, for an accepted program, they produce
the same observations: the same output, the same DOM writes, the same calls into JavaScript with the
same arguments marshalled by value (\S\ref{sec:soundness}).

\paragraph{Scope.} Records, tuples and constructors are not written in place by this checker,
although the machinery extends to them (\S\ref{sec:compiler-records}). The checker does not decide
whether an in-place update is \emph{worthwhile}: every accepted update is in place, and there is no
cost model.
